% SGA 3, Expose IV, English reconstruction.
% Scope: Lemma 4.4.7, completion of the proof of Proposition 4.4.3, and
% Remarks 4.4.8 and 4.4.8 bis, stopping before editor's note 38 and Lemma
% 4.4.8.1.
% Authority: Polo--Gille Exp4-13oct24.pdf, component-local pages 40--41
% (both partial; combined-reader pages 228--229).
% The external English Markdown was used only as a comparison.

\setcounter{footnote}{34}

\medskip
\noindent\SGARefTarget{sga3:iv:lemma:4:4:7}\textbf{引理 4.4.7.}
\textit{层的每个态射族 \(\{F_i\to F\}\) 都分解为一个覆盖族
\(\{F_i\to G\}\) 后接一个单态射 \(G\to F\).}

事实上, 只须取 \(G\) 为所给族的层像.

\medskip
\noindent\textit{命题 \SGARefLink{sga3:iv:proposition:4:4:3}{4.4.3} 的证明.}
已在 \SGARefLink{sga3:iv:lemma:4:4:6}{4.4.6} 中看到,
\SGARefLink{sga3:iv:proposition:4:4:3:item:iii:01}{(iii)} 蕴含
\SGARefLink{sga3:iv:proposition:4:4:3:item:ii:01}{(ii)}, 而
\SGARefLink{sga3:iv:proposition:4:4:3:item:ii:01}{(ii)} 显然蕴含
\SGARefLink{sga3:iv:proposition:4:4:3:item:i:01}{(i)}.
最后, 设 \(\{F_i\to F\}\) 是满态射族.
由引理~\SGARefLink{sga3:iv:lemma:4:4:7}{4.4.7},
它分解为一个覆盖族后接一个单态射. 后一个单态射被一个满态射族支配,
\footnote{编者注: 回忆 (参见编者注 17), 若态射族 \(F_i\to H\)
经过态射 \(g:G\to H\) 分解, 则称 \(g\) \emph{被}此族\emph{支配}.}
因此它本身也是满态射; 因而由 \SGARefLink{sga3:iv:lemma:4:4:4}{4.4.4}, 它是同构.

\medskip
\noindent\SGARefTarget{sga3:iv:remark:4:4:8}\textbf{注 4.4.8.}
\footnote{编者注: 原文有两段都编号为 \SGARefLink{sga3:iv:lemma:4:4:7}{4.4.7},
现已改正编号, 以与后文引用相符.}
由于族 \(\mathscr F=\{F_i\to F\}\) 的预层像分离,
由相伴层的构造可知, 命题~\SGARefLink{sga3:iv:proposition:4:4:3}{4.4.3}
中的条件还与下列条件等价:
\begin{enumerate}
  \item[(v)]\SGARefTarget{sga3:iv:remark:4:4:8:item:v:01}
    对每个 \(S\in\operatorname{Ob}\mathscr C\),
    每个 \(f\in F(S)\) 局部地属于 \(\mathscr F\) 的像; 即,

  \item[(vi)]\SGARefTarget{sga3:iv:remark:4:4:8:item:vi:01}
    对每个 \(S\in\operatorname{Ob}\mathscr C\) 和每个 \(f\in F(S)\),
    使 \(f\) 在 \(F(S')\) 中的像属于某个 \(F_i(S')\) 的像的
    态射 \(S'\to S\) 所成的集合, 是 \(S\) 的加细;

  \item[(vii)]\SGARefTarget{sga3:iv:remark:4:4:8:item:vii:01}
    若拓扑由一个预拓扑 \(S\mapsto\mathscr R(S)\) 定义,
    则对每个 \(S\in\operatorname{Ob}\mathscr C\) 和每个 \(f\in F(S)\),
    存在族 \(\{S_j\to S\}\in\mathscr R(S)\),
    使得对每个 \(j\), \(f\) 在 \(F(S_j)\) 中的像 \(f_j\)
    属于某个 \(F_i(S_j)\) 的像.
\end{enumerate}

\medskip
\noindent\SGARefTarget{sga3:iv:remark:4:4:8:bis}\textbf{注 4.4.8 bis.}
若层 \(F\) 可表, 则前述条件还等价于:
\begin{enumerate}
  \item[(viii)]\SGARefTarget{sga3:iv:remark:4:4:8:bis:item:viii:01}
    所有满足下述条件的态射 \(T\to F\) (其中
    \(T\in\operatorname{Ob}\mathscr C\)) 所成的集合是 \(F\) 的加细:
    存在一个 \(i\) 及一个交换图
    \[
    \begin{array}{ccc}
       &F_i&\\[-0.4ex]
       &\big\downarrow&\\[-0.4ex]
      T&\longrightarrow&F
    \end{array}
    \]
\end{enumerate}

事实上, 若 \SGARefLink{sga3:iv:remark:4:4:8:bis:item:viii:01}{(viii)} 成立,
则各 \(F_i\) 的预层像被 \(F\) 的一个加细支配,
\footnote{编者注: \textquotedblleft 包含\textquotedblright{}
已改为 \textquotedblleft 被支配\textquotedblright; 参见编者注 17.}
这蕴含此族为覆盖族. 反过来, 对 \(\operatorname{id}_F\in F(F)\)
应用 \SGARefLink{sga3:iv:remark:4:4:8:item:vi:01}{(vi)} 即可.

\SGARefTarget{sga3:iv:remark:4:4:8:bis:item:b:01}
形象地说, \SGARefLink{sga3:iv:remark:4:4:8:bis:item:viii:01}{此条件}意味着:
\textit{在 \(F\) 上局部地, 存在一个 \(i\), 使 \(F_i\to F\) 有截面}.
特别地, 若态射
\[
  G\longrightarrow F,
\]
中 \(G\) 是层, \(F\) 是可表层, 则此态射是覆盖态射,
当且仅当它在 \(F\) 上局部地有截面.

% Translation coverage: entire assigned source component.
% Source lines: 85; translated body lines before this audit comment: 89.
% Math preserved: 42 inline formulas, 2 display blocks.
% Reference targets: 8/8; links: 13/13; labels: 0/0; footnotes: 3/3.
% Unresolved TODOs: 0; compile and visual review delegated to the volume integrator.
